\section{Methodology}\label{sec:method}

We propose \textbf{VOILA} (\underline{V}alue-\underline{O}f-\underline{I}nformation guided \underline{L}atency-\underline{A}ware fidelity selection), a query-adaptive framework for selecting the minimum-cost input fidelity required to answer a multimodal query accurately. VOILA operates entirely at inference time and is model-agnostic, requiring no modification to the underlying vision--language model (VLM).

The key challenge VOILA addresses is \emph{pre-retrieval decision-making}: \textbf{fidelity choices must be made \emph{before} executing expensive visual retrieval or VLM inference}.
As a result, VOILA must reason about the expected utility of acquiring additional information using only signals available \emph{prior} to retrieval.

\subsection{Problem Setup}
\label{sec:problem_setup}

Let $(q, x)$ denote a visual question $q$ paired with an image $x$.
We assume access to a discrete, ordered set of input fidelities
\[
\mathcal{F} = \{f_1, f_2, \dots, f_K\},
\]
sorted by increasing acquisition cost.
Lower fidelities correspond to inexpensive representations (e.g., captions or thumbnails),
while higher fidelities provide richer visual information (e.g., full-resolution images).

Each fidelity $f \in \mathcal{F}$ incurs a cost $c(f) > 0$, capturing retrieval latency,
bandwidth, and compute overhead.
Given a fidelity $f$, the vision--language model (VLM) would produce an answer
$\hat{a}_f$; however, executing the VLM at multiple fidelities is prohibitively expensive at inference time.

\textbf{Pre-retrieval constraint.}
VOILA operates in a \emph{pre-retrieval} regime:
for each query, the system must select a fidelity \textbf{\emph{before}} (i) retrieving the
corresponding visual representation and (ii) executing the VLM.
As a result, fidelity selection cannot rely on post-hoc model signals
(e.g., token probabilities or hidden activations) and must instead be made using
question-only information.

Our objective is to select the fidelity that maximizes expected utility,
\[
f^\star
=
\arg\max_{f \in \mathcal{F}}
\left(
\Pr(\hat{a}_f \text{ is correct} \mid q) - \lambda c(f)
\right),
\]
where the probability of correctness must be estimated without retrieving fidelity $f$
or running the VLM at that fidelity.



\subsection{Question-Conditioned Success Scoring}

To estimate fidelity-specific success probabilities in a pre-retrieval setting, we learn
question-conditioned predictors using supervised data.
For each fidelity $f$, we construct a training set by executing the VLM at fidelity $f$
on labeled examples and recording whether the produced answer is correct.

Let $y_f \in \{0,1\}$ denote the correctness label for fidelity $f$.
We extract a question feature representation $\phi(q)$ consisting of:
\begin{itemize}
    \item TF--IDF features of the question text,
    \item lexical statistics (e.g., length, numeric indicators),
    \item coarse question-type features (e.g., counting, color, spatial).
\end{itemize}

Using these features, we train a regressor
\[
r_f(q) = g_f(\phi(q)),
\]
implemented as a Gradient Boosting Regressor (GBR), to produce a raw score that correlates
with the likelihood that the VLM will answer correctly at fidelity $f$.
GBR is used as a robust default for sparse, mixed feature spaces, though the framework is
agnostic to the specific regression model.



\subsection{Fidelity-Specific Probability Calibration}

The raw scores $r_f(q)$ produced by the regressor are not guaranteed to be calibrated
probabilities.
To obtain reliable probability estimates suitable for decision-making, we apply
fidelity-specific post-hoc calibration.

For each fidelity $f$, we fit an isotonic regression model
\[
\psi_f: r_f(q) \mapsto \hat{p}_f(q),
\]
on a held-out calibration split, where $\hat{p}_f(q) \in [0,1]$ estimates the probability
that the VLM would answer correctly at fidelity $f$.
Isotonic regression enforces monotonicity while remaining non-parametric, allowing it to
correct systematic over- and under-confidence without assuming a particular functional
form.

At inference time, only the calibrated predictor $\hat{p}_f(q)$ is used; no additional
calibration or model execution is required.

\subsection{Value-of-Information–Based Fidelity Selection}

Given calibrated success probabilities $\hat{p}_f(q)$, VOILA selects fidelities by
comparing their expected utility.
Starting from the lowest-cost fidelity, the system evaluates whether escalating to a
higher fidelity is beneficial by computing the Value of Information (VOI):
\[
\mathrm{VOI}(f_j \mid f_i)
=
\hat{p}_{f_j}(q) - \hat{p}_{f_i}(q) - \lambda c(f_j),
\]
where $\lambda$ controls the trade-off between accuracy and cost.

VOILA incrementally increases fidelity only when the expected improvement in correctness
outweighs the additional acquisition cost.

\begin{algorithm}[t]
\caption{VOILA Inference Procedure}
\label{alg:voila}
\begin{algorithmic}[1]
\STATE Initialize $f \leftarrow f_1$
\STATE Compute $\hat{p}_f(q)$
\FOR{$f' \in \mathcal{F}$ in ascending cost}
    \STATE Compute $\hat{p}_{f'}(q)$
    \STATE $\mathrm{VOI} \leftarrow \hat{p}_{f'}(q) - \hat{p}_f(q) - \lambda c(f')$
    \IF{$\mathrm{VOI} > \tau$}
        \STATE $f \leftarrow f'$
        \STATE $\hat{p}_f(q) \leftarrow \hat{p}_{f'}(q)$
    \ELSE
        \STATE \textbf{break}
    \ENDIF
\ENDFOR
\STATE \textbf{return} answer produced at fidelity $f$
\end{algorithmic}
\end{algorithm}
